\subsection{最初的推论}\label{sec:18j}
\stref{命题}{18.23}最直接的应用，是将 Hurwitz 移动与变异联系起来。取消失路径的一个特选基
\SourcePage{295}
$\symbfit{\gamma}=(\gamma_1,\ldots,\gamma_m)$，考虑相应的 (B) 型等变膜集 $\widetilde{\symbfit{\Delta}}_{\symbfit{\gamma}}^\#$，其中分次任意选取；同时考虑带诱导膜结构的相应消失圈集 $\symbfit V_{\symbfit{\gamma}}^\#$。由\stref{命题}{18.14}，已知 $\widetilde{\symbfit{\Delta}}_{\symbfit{\gamma}}^\#$ 是例外集。\stref{命题}{18.23}说明，若对特选基施加\eref{16.7}中定义的一次初等 Hurwitz 移动，所得结果就是\eref{5.24}意义下该例外集的一次初等变异。选择相关记号时，我们已隐含地承认了这一对应，因此可写成
\begin{equation}\label{eq:18.29}
\widetilde{\symbfit{\Delta}}_{\mathcal L_k\symbfit{\gamma}}^\#
\cong\mathcal L_k\widetilde{\symbfit{\Delta}}_{\symbfit{\gamma}}^\#.
\end{equation}
严格地说，为使\eref{18.29}成为 $D\F(\pi)$ 中例外集的真正同构，分次必须按\stref{命题}{18.23}的方式专门选择；否则，它只在相差平移的意义下是同构，而这已完全足够。显然，只需改变所涉及对象集的记号，同样的结论就适用于逆移动 $\mathcal L_k^{-1}$，再通过复合，适用于任意变异。换言之，辫群 $\operatorname{Br}_m$ 在模去平移与同构的例外集上的代数作用，见\secref{5j}，可与它在消失路径特选基上的几何作用，见\secref{16d}，等同起来。

固定一个 $\symbfit{\gamma}$。暂以 $\mathcal T$ 表示 $\Tw\F(\pi)$ 中由对象 $\widetilde\Delta_{\gamma_1}^\#,\ldots,\widetilde\Delta_{\gamma_m}^\#$ 生成的三角子范畴。依定义，$\widetilde{\symbfit{\Delta}}_{\symbfit{\gamma}}^\#$ 是该子范畴中的满例外集，所有变异之后也仍如此。由\eref{18.29}，这意味着，若对基施加 Hurwitz 移动，再考虑与所得基相应的对象集，则其全部对象仍在 $\mathcal T$ 中，并生成该范畴。由于 Hurwitz 移动在特选基上可迁地作用，$\mathcal T$ 与 $\symbfit{\gamma}$ 无关。事实上，\stref{命题}{18.17}说明，任意 (U) 型等变膜都属于 $\mathcal T$，因此，由 $\F(\pi)$ 的定义，必有 $\mathcal T=\Tw\F(\pi)$。另一方面，\stref{引理}{5.22}说明，$\mathcal T$ 拟等价于与 $\widetilde{\symbfit{\Delta}}_{\symbfit{\gamma}}^\#$ 相应的有向 $\Ai$ 子范畴上的扭曲复形范畴；再由\stref{命题}{18.14}，该有向子范畴拟同构于 $\F^\to(\symbfit V_{\symbfit{\gamma}}^\#)$。\znote{287}{定理的证明是四步，第一步最容易在读时被跳过：(1) 引理 18.2 与命题 18.23 给出\eref{18.29}，即初等 Hurwitz 移动对应初等变异，于是「满例外集的变异」与「基的 Hurwitz 移动」同步；(2) Hurwitz 移动在特选基集合上可迁地作用，故 $\mathcal T$ 与所取基无关；(3) 命题 18.17 把任意 (U) 型膜收进 $\mathcal T$，而 $\F(\pi)$ 由 $\widetilde\Delta_{k}^\#$ 与这些膜生成，故 $\mathcal T=\Tw\F(\pi)$；(4) 引理 5.22 把 $\mathcal T$ 实现为 $\widetilde{\symbfit{\Delta}}_{\symbfit{\gamma}}^\#$ 上扭曲复形范畴，命题 18.14 再把后者换成 $\Tw\F^\to(\symbfit V_{\symbfit{\gamma}}^\#)$。注意定理本身不另列假设：它是在 §18f 起固定的 $\operatorname{char}\K\ne2$ 之下陈述的，代价在那里一次付清；四步中只有 (1) 真正用到二重分歧覆盖。[来源：原书 §18j（印刷页 287）]}因此：

\StatementNumber{24}\begin{theorem}\label{st:18.24}
$\Tw\F^\to(\symbfit V_{\symbfit{\gamma}}^\#)$ 拟等价于 $\Tw\F(\pi)$。
\end{theorem}

一个直接推论是得到\stref{推论}{17.21}的另一证明；事实上，也得到\stref{定理}{17.20}的另一证明，因为任意 Lagrangian 球面集都可实现为某个适当构造的 Lefschetz 纤维化的消失圈基。这比原来的证明有所改进：通过将 $\Tw\F^\to(\symbfit V_{\symbfit{\gamma}}^\#)$ 与显然不涉及基的选择的 $\Tw\F(\pi)$ 等同起来，证明了它与 $\symbfit{\gamma}$ 无关。一个具体收获是，可以为不同 $\symbfit{\gamma}$ 对应的范畴 $\Tw\F^\to(\symbfit V_{\symbfit{\gamma}}^\#)$ 构造典范的拟等价，至少在相应函子范畴 $H^0(\operatorname{fun})$ 中的同构意义下如此。
\SourcePage{296}
为此付出的唯一代价，是技术假设 $\operatorname{char}(\K)\ne2$；使用二重分歧覆盖这一技巧，迫使我们作出该假设。

其次，注意\stref{命题}{18.13}与\stref{引理}{3.23}给出一个在上同调上全忠实的函子 $\Tw\F(E)\to\Tw\F(\pi)$。结合\stref{定理}{18.24}，得到：
\StatementNumber{25}\begin{corollary}\label{st:18.25}
对任意 $\symbfit{\gamma}$，存在 $\F(E)$ 到 $\Tw\F^\to(\symbfit V_{\symbfit{\gamma}}^\#)$ 的在上同调上全忠实的嵌入。
\end{corollary}

关键在于，有向 Fukaya 范畴是在 $\F(M)$ 内部构造的，因此，从抽象角度说，总空间 $E$ 中的 Floer 上同调计算可以完全归约为纤维中的计算。下面尝试把这些讨论变得更加具体。\znote{288}{注意这里是「拟满忠实」而非「等价」：$\F(E)\hookrightarrow\Tw\F^\to(\symbfit V_{\symbfit{\gamma}}^\#)$ 只保证 $\Hom$ 上的同构，并不保证每个对象都来自 $E$ 中的膜，故 $E$ 的 Fukaya 范畴实现为消失圈所给有向 $\Ai$ 范畴的扭曲复形范畴的一个全子范畴。这句话的实际含义因此是：全空间的 Floer 数据被纤维内一列球面（连同它们的相交数与高阶运算）完全承载，这正是 Picard--Lefschetz 理论把整体问题化为组合数据的着力点。也正因如此，命题 18.13 的证明分两步：先说明 $E^{\mathrm{shrink}}$ 内的膜已足以算出 $\F(E)$（用定理 2.9 反演），再说明它们与 $\widetilde{\A}^{\mathrm{inv}}$ 的一个全子范畴一致。[来源：原书 §18j、命题 18.13、定理 18.24（印刷页 282、287、288）]}
